\chapter{Arithmetic on a clock}
\label{chap:modular}
\chaptertag{Extended Euclid, Chinese remainders, and interpolation}

\begin{bigidea}
The extended gcd algorithm is the computational engine of modular arithmetic.
It tells us when division is possible and supplies the inverse when it is.
We will use those inverses to reconstruct an integer from its remainders,
then recognize the same construction in polynomial interpolation.
Finally, a reversible recurrence will explain why Fibonacci residues repeat.
\end{bigidea}

\begin{definition}
For $n\ge2$, we write $a\congr b\pmod n$ if $n\divides(a-b)$. The equivalence class of $a$ is written $[a]$ or $\bar a$.
\end{definition}

Congruence remembers the remainder and forgets the number of complete turns around the clock.

\begin{center}
\begin{tikzpicture}[scale=1.0]
  \draw[navy!35,thick] (0,0) circle (1.65);
  \foreach \k in {0,...,11}{
    \coordinate (P\k) at ({90-30*\k}:1.65);
    \fill[teal] (P\k) circle (2pt);
    \node at ({90-30*\k}:1.94) {$\k$};
  }
  \draw[-{Stealth[length=3mm]},coral,line width=1.5pt] (P9) arc[start angle=180,end angle=30,radius=1.65];
  \node[coral,align=left,anchor=west] at (2.6,.2) {$9+5\congr2\pmod{12}$\\one complete turn disappears};
\end{tikzpicture}
\end{center}

\section{Legal arithmetic}

If $a\congr b\pmod n$ and $c\congr d\pmod n$, then
\[
a+c\congr b+d\pmod n,\qquad ac\congr bd\pmod n.
\]
Powers are therefore legal too. Division is different.

These rules do not depend on the chosen representatives. For example, if
$a=b+kn$ and $c=d+\ell n$, then
$ac-bd=n(b\ell+dk+k\ell n)$. Thus addition and multiplication are
well-defined operations on the residue classes $\Z/n\Z$. Reduce after
each step or at the end: the resulting residue class is the same.

\section{Extended gcd: division comes with a certificate}
\label{sec:modular-egcd}

\begin{definition}
A residue $a$ is a \emph{unit modulo $n$} if there exists $b$ with $ab\congr1\pmod n$. Such a $b$ is an inverse of $a$.
\end{definition}

\begin{theorem}[Invertibility criterion]
$a$ has an inverse modulo $n$ exactly when $\gcd(a,n)=1$.
\end{theorem}

\begin{proof}
The congruence $ab\congr1\pmod n$ means $ab+nk=1$ for some $k$, so any
common divisor of $a,n$ divides $1$. Conversely, when $\gcd(a,n)=1$,
extended Euclid produces $u,v$ such that $au+nv=1$. Reducing this
\emph{integer identity} modulo $n$ gives $au\congr1$: the coefficient
$u$, reduced modulo $n$, is the required inverse.
\end{proof}

The theorem is an algorithm, not a direction to try every residue until
one works. Here is a forward version of the backward substitutions from
Chapter~1. First reduce $a$ so that $0\le a<n$. Keep three entries in each row:
\[
 r_i=ns_i+at_i.
\]
Start with $(r_0,s_0,t_0)=(n,1,0)$ and $(r_1,s_1,t_1)=(a,0,1)$.
Whenever $r_i\ne0$, divide $r_{i-1}=q_i r_i+r_{i+1}$ and update
\[
 (r_{i+1},s_{i+1},t_{i+1})
 =(r_{i-1},s_{i-1},t_{i-1})-q_i(r_i,s_i,t_i).
\]
The same subtraction in all three columns preserves the displayed
identity. The last nonzero row gives the gcd and its B\'ezout coefficients.

\begin{example}[An inverse without guessing]
Find $17^{-1}\pmod{43}$. The divisions are
\[
43=2\cdot17+9,\quad17=9+8,\quad9=8+1,\quad8=8\cdot1.
\]
The coefficient table records $r=43s+17t$ throughout:
\begin{center}
\begin{tabular}{rrrl}
\toprule
$r$ & $s$ & $t$ & Certificate \\
\midrule
43 & 1 & 0 & $43=1\cdot43+0\cdot17$ \\
17 & 0 & 1 & $17=0\cdot43+1\cdot17$ \\
9 & 1 & $-2$ & $9=43-2\cdot17$ \\
8 & $-1$ & 3 & $8=-43+3\cdot17$ \\
1 & 2 & $-5$ & $1=2\cdot43-5\cdot17$ \\
0 & $-17$ & 43 & $0=-17\cdot43+43\cdot17$ \\
\bottomrule
\end{tabular}
\end{center}
Hence $17^{-1}\congr-5\congr38\pmod{43}$. Check the result directly:
$17\cdot38=646=15\cdot43+1$.
\end{example}

\begin{checkpoint}
In a certificate $au+nv=1$, which coefficient is the inverse of $a$
modulo $n$? Read the coefficient of $a$, even if the table began with
$n$ in its first row. The final multiplication check catches a swapped
coefficient or a sign error.
\end{checkpoint}

\begin{warningbox}
Cancellation can fail. Modulo $6$, $2\cdot1\congr2\cdot4$, but $1\not\congr4$. You may cancel $a$ modulo $n$ only when $a$ is a unit, or after carefully reducing the modulus.
\end{warningbox}

For a prime $p$, every nonzero residue is a unit. We write
$\F_p=\Z/p\Z$ and call it a \emph{field}: addition and multiplication obey
the usual commutative arithmetic laws, and every nonzero element permits
division. The rational and real numbers are other examples of fields.
For composite $n$, $\Z/n\Z$ need not permit division by a nonzero residue.

\section{Solving a linear congruence with the same algorithm}

\begin{proposition}[Linear congruence]
The congruence $ax\congr b\pmod n$ has a solution exactly when $d=\gcd(a,n)$ divides $b$. When it does, there are exactly $d$ incongruent solutions modulo $n$.
\end{proposition}

\begin{proof}[Construction and completeness]
Run extended Euclid once to obtain $au+nv=d$. If $ax\congr b\pmod n$,
then $b=ax-nk$ for an integer $k$, so $d\divides b$ is necessary.
If $d\divides b$, multiply the certificate by $b/d$ to obtain a solution
\[
 x_0=u\frac bd.
\]
If $d=n$, then both $a$ and $b$ are multiples of $n$, so all $n$ residues
are solutions. Otherwise $n/d\ge2$, and any two solutions satisfy
$(a/d)(x-x_0)\congr0\pmod{n/d}$. Since $a/d$ is a unit modulo $n/d$,
this is equivalent to $x\congr x_0\pmod{n/d}$. Therefore the complete list
modulo $n$ is
\[
 x_0,\ x_0+\frac nd,\ \ldots,\ x_0+(d-1)\frac nd.
\]
These $d$ residues are distinct.
\end{proof}

\begin{example}[A nonunit can give several solutions]
For $14x\congr21\pmod{35}$, extended Euclid gives
$7=35-2\cdot14$. Multiplying by $3$ gives $21=3\cdot35-6\cdot14$,
so $x\congr-6\congr4\pmod5$. The seven solutions modulo $35$ are
$4,9,14,19,24,29,34$. The same calculation rules out
$14x\congr20\pmod{35}$, because $7\ndivides20$.
\end{example}

\section{Reassembling remainders}
\label{sec:crt}

Suppose a number leaves remainder $2$ modulo $3$ and remainder $3$ modulo $5$. Listing possibilities gives $8$ modulo $15$. The Chinese remainder theorem says this always works for coprime clocks.

\begin{theorem}[Chinese remainder theorem]
If $n_1,\dots,n_r$ are pairwise coprime, then for any $a_1,\dots,a_r$, the system
\[
x\congr a_i\pmod{n_i}\qquad(i=1,\dots,r)
\]
has a solution, unique modulo $N=n_1\cdots n_r$.
\end{theorem}

\begin{proof}[Construction]
Set $N_i=N/n_i$. Since $\gcd(N_i,n_i)=1$, \emph{use extended Euclid} to
find $M_i,v_i$ with $N_iM_i+n_iv_i=1$. Then
\[
x=\sum_{i=1}^r a_iN_iM_i
\]
has every prescribed remainder: modulo $n_i$, its $i$th term is $a_i$ and every other term is $0$. If $x,y$ both work, each $n_i$ divides $x-y$; coprimality makes $N\divides x-y$.
\end{proof}

The quantities $e_i=N_iM_i$ are \emph{selectors}: $e_i$ is $1$ at its
own modulus and $0$ at every other modulus. The term $a_i e_i$ installs
one prescribed remainder without disturbing the others.

\begin{example}[Construct, reduce, and check]
Solve $x\congr2\pmod3$, $x\congr3\pmod5$, $x\congr2\pmod7$.
Here $N=105$, and the construction is
\begin{center}
\begin{tabular}{rrrrr}
\toprule
$n_i$ & $a_i$ & $N_i$ & $M_i$ & $e_i=N_iM_i$ \\
\midrule
3 & 2 & 35 & 2 & 70 \\
5 & 3 & 21 & 1 & 21 \\
7 & 2 & 15 & 1 & 15 \\
\bottomrule
\end{tabular}
\end{center}
For instance, $2\cdot35-23\cdot3=1$ certifies the first inverse.
Thus $x=2\cdot70+3\cdot21+2\cdot15=233\congr23\pmod{105}$.
Finally, $23$ has remainders $2,3,2$, as required. Every solution is
$23+105k$ for an integer $k$.
\end{example}

\begin{center}
\begin{tikzpicture}[node distance=8mm]
\node[draw=teal,fill=mist,rounded corners,minimum width=21mm] (z15) {$\Z/15\Z$};
\node[draw=navy!55,fill=white,rounded corners,below left=of z15,minimum width=21mm] (z3) {$\Z/3\Z$};
\node[draw=navy!55,fill=white,rounded corners,below right=of z15,minimum width=21mm] (z5) {$\Z/5\Z$};
\draw[-{Stealth},thick,teal] (z15)--(z3) node[midway,left,font=\scriptsize] {$x\mapsto x\bmod3$};
\draw[-{Stealth},thick,teal] (z15)--(z5) node[midway,right,font=\scriptsize] {$x\mapsto x\bmod5$};
\node[below=12mm of z15,muted] {One clock of size $15$ is equivalent to two independent clocks.};
\end{tikzpicture}
\end{center}

The map from one clock to its component clocks preserves addition and
multiplication, since reduction does. CRT says it is also a bijection.
For example, the selectors modulo $15$ are $10$ and $6$, giving
$x\congr10a+6b\pmod{15}$ for remainders $a$ modulo $3$ and $b$ modulo $5$.

\begin{proposition}[When two clocks share a factor]
The conditions $x\congr a\pmod m$, $x\congr b\pmod n$ are compatible
exactly when $a\congr b\pmod{\gcd(m,n)}$. If compatible, the solution is
unique modulo $\operatorname{lcm}(m,n)$.
\end{proposition}
\begin{proof}
Write $x=a+mt$. The remaining condition is $mt\congr b-a\pmod n$,
so the linear-congruence construction supplies a solution exactly when
$\gcd(m,n)\divides b-a$. Differences between two solutions are precisely
the common multiples of $m$ and $n$.
\end{proof}

\section{Chinese remainders become Lagrange interpolation}
\label{sec:crt-interpolation}

CRT builds an integer with prescribed remainders. Interpolation builds
a polynomial with prescribed values. Polynomial division makes these
the same kind of problem.

Let $K$ be a field, such as $\Q$, $\R$, or $\F_p$. Write $K[X]$ for
polynomials in an indeterminate $X$ with coefficients in $K$. Division
by $X-c$ gives
\[
 f(X)=(X-c)q(X)+r,
\]
where $r$ is constant. Substituting $X=c$ shows $r=f(c)$. In other words,
\[
 f(X)\congr b\pmod{X-c}\quad\Longleftrightarrow\quad f(c)=b.
\]
Here polynomial congruence means that the difference is divisible by the
polynomial modulus.

\subsection*{The same construction, with polynomial moduli}

The division algorithm works in $K[X]$: cancel the leading term using
division in $K$, subtract, and continue while the degree decreases.
Extended Euclid therefore works too. For coprime polynomials $A,B$,
it produces $U,V$ with $UA+VB=1$.

Consequently, for pairwise coprime monic moduli $m_1,\ldots,m_r$, set
$P=\prod_i m_i$ and $P_i=P/m_i$, and use extended Euclid to find
$U_iP_i+V_im_i=1$. The polynomial
\[
 f=\sum_i b_iP_iU_i
\]
has prescribed polynomial remainders $b_i$ modulo $m_i$.
The proof is the integer CRT proof: reduce at each modulus separately.
Uniqueness is modulo $P$, since the pairwise coprime factors $m_i$ all
divide the difference of two solutions. Dividing by the monic polynomial
$P$ selects one representative of degree less than $\deg P$.

For interpolation, take $m_i=X-c_i$ with the $c_i$ distinct. Two such
linear polynomials are coprime: subtracting them gives the nonzero
constant $c_j-c_i$, which is invertible in $K$.
In this special case, the inverse required by CRT is particularly simple.
For
\[
 P_i(X)=\prod_{j\ne i}(X-c_j),
\]
reduction modulo $X-c_i$ gives the nonzero number
$P_i(c_i)=\prod_{j\ne i}(c_i-c_j)$.
Thus define
\[
 \ell_i(X)=\frac{P_i(X)}{P_i(c_i)}
 =\prod_{j\ne i}\frac{X-c_j}{c_i-c_j}.
\]
These are polynomial selectors: $\ell_i(c_i)=1$ and
$\ell_i(c_j)=0$ for $j\ne i$.

\begin{theorem}[Lagrange interpolation]
Let $c_1,\ldots,c_r$ be distinct points of a field $K$.
Given values $b_i\in K$, the unique polynomial of degree less than $r$
with $f(c_i)=b_i$ for every $i$ is
\[
 f(X)=\sum_{i=1}^r b_i\ell_i(X).
\]
Every polynomial with these values has the form
$f(X)+h(X)\prod_i(X-c_i)$, with $h\in K[X]$.
\end{theorem}

\begin{proof}
The selector values prove existence. If $f,g$ have the required values,
each $X-c_i$ divides $f-g$, and hence so does their product. A nonzero
multiple of that degree-$r$ product cannot have degree less than $r$.
This proves uniqueness and the description of all solutions.
The same argument proves the useful \emph{root bound}: a nonzero
polynomial of degree $d$ over a field has at most $d$ distinct roots.
\end{proof}

\begin{center}
\small
\begin{tabularx}{\textwidth}{p{.18\textwidth}XX}
\toprule
 & Integer CRT & Lagrange interpolation \\
\midrule
Local data & $x\congr a_i\pmod{n_i}$ & $f(c_i)=b_i$ \\
Moduli & pairwise coprime $n_i$ & distinct linear factors $X-c_i$ \\
Selector & $N_i(N_i^{-1}\bmod n_i)$ & $P_i(X)/P_i(c_i)$ \\
Formula & $\sum_i a_i e_i$ & $\sum_i b_i\ell_i(X)$ \\
Ambiguity & multiples of $N$ & multiples of $\prod_i(X-c_i)$ \\
Choice & $0\le x<N$ & $\deg f<r$ \\
\bottomrule
\end{tabularx}
\end{center}

\begin{example}[Three values determine a quadratic]
Over $\Q$, prescribe $f(0)=1$, $f(1)=2$, $f(2)=5$. The selectors are
\[
 \ell_0=\frac{(X-1)(X-2)}2,\qquad
 \ell_1=-X(X-2),\qquad
 \ell_2=\frac{X(X-1)}2.
\]
Combining them gives $f=\ell_0+2\ell_1+5\ell_2=X^2+1$.
All other solutions are $X^2+1+X(X-1)(X-2)h(X)$.

The same computation works in $\F_5$, where the prescribed values are
$1,2,0$ and division by $2$ means multiplication by $3$.
The certificate $2\cdot3-5=1$ makes that division legitimate.
The result is again $X^2+1$, now with coefficients modulo $5$.
\end{example}

\begin{warningbox}
Distinct interpolation points must have invertible differences.
This is automatic in a field. It can fail for composite moduli:
over $\Z/8\Z$, the distinct points $0$ and $2$ have difference $2$,
which has no inverse. Indeed, $f(2)-f(0)$ is even for every polynomial,
so the data $f(0)=0$, $f(2)=1$ cannot be interpolated.
\end{warningbox}

\begin{checkpoint}
Both reconstruction formulas first make an object that vanishes at every
unwanted coordinate, then multiply by an inverse to make its remaining
coordinate equal to $1$. Identify those two steps before memorizing either
formula.
\end{checkpoint}


\section{Fast powers and patterns}

Modular reduction prevents numbers from exploding. To compute $7^{222}\pmod{13}$, square and reduce:
\[
7^2\congr10,\quad 7^4\congr9,\quad7^8\congr3,\quad7^{16}\congr9\pmod{13},
\]
then use the binary expansion $222=128+64+16+8+4+2$. Later, Fermat's theorem will make this even faster.

\begin{labbox}
Make a multiplication table modulo $n$ and highlight the rows that are permutations of all residues. Conjecture which row labels have inverses. Then compare with the gcd criterion.
\end{labbox}

\section{Fibonacci residues: why the pattern must repeat}
\label{sec:fibonacci-periodicity}

Set $F_0=0$, $F_1=1$, and $F_{n+2}=F_{n+1}+F_n$. Modulo an integer
$m\ge2$, the next term depends on the previous \emph{two} terms.
The appropriate state is therefore
\[
 S_n=(F_n,F_{n+1})\pmod m,
 \qquad T(a,b)=(b,a+b)\pmod m.
\]
There are only $m^2$ pairs, so a state must eventually repeat.
But finiteness alone only gives eventual periodicity. Here we can do better:
\[
 T^{-1}(u,v)=(v-u,u)\pmod m.
\]
Every step has a unique predecessor, using only subtraction.

\begin{theorem}[Fibonacci periodicity from the starting pair]
For every $m\ge2$, the Fibonacci sequence modulo $m$ is periodic from
the start. Its least period $\pi(m)$ is the smallest positive $t$ for which
\[
 (F_t,F_{t+1})\congr(0,1)\pmod m.
\]
In particular, $\pi(m)\le m^2-1$.
\end{theorem}

\begin{proof}
If $S_i=S_j$ with $i<j$, apply the inverse step $i$ times to get
$S_0=S_{j-i}$. Repeating the forward step now gives a period beginning
at $S_0$, not just after an initial segment. Conversely, a period of the
sequence must return both $F_0$ and $F_1$, so it returns the starting pair.
Before the first return all states are distinct. The pair $(0,0)$ is fixed
by $T$ and has no other predecessor, so the orbit of $(0,1)$ never visits
it. There are at most $m^2-1$ available states.
\end{proof}

\begin{example}[A zero is not yet a restart]
Modulo $3$, the successive states are
\[
\begin{array}{c|rrrrrrrrr}
n&0&1&2&3&4&5&6&7&8\\\hline
F_n&0&1&1&2&0&2&2&1&0\\
F_{n+1}&1&1&2&0&2&2&1&0&1
\end{array}
\]
The pair first returns at $n=8$, so $\pi(3)=8$.
Although $F_4\congr0$, its successor is $2$, not $1$.
Looking only for a repeated Fibonacci value would give the wrong answer.
\end{example}

\begin{proposition}[CRT combines periods]
If $\gcd(m,n)=1$, then
\[
 \pi(mn)=\operatorname{lcm}(\pi(m),\pi(n)).
\]
\end{proposition}

\begin{proof}
CRT says that a pair returns modulo $mn$ exactly when it returns modulo
both $m$ and $n$. The return times modulo $m$ are precisely the multiples
of $\pi(m)$: divide any return time by $\pi(m)$ and use the minimality of
the remainder. Likewise for $n$. The first common return time is their
least common multiple.
\end{proof}

For example, the orbit modulo $2$ has length $3$, so
$\pi(6)=\operatorname{lcm}(3,8)=24$.
The connection is operational: CRT reassembles the two coordinates of
every state as well as the arithmetic on the clocks.

\begin{labbox}
For a small modulus, start at $(0,1)$, repeatedly replace $(a,b)$ by
$(b,a+b\bmod m)$, and stop on the first return to $(0,1)$.
The theorem proves this procedure terminates and that the answer is exact.
Record pairs, not just zeros. These small experiments need no extension fields.
\end{labbox}

\subsection*{What we can prove now, and what comes next}

The reversible-pair argument proves periodicity for \emph{every} modulus,
prime or composite. The estimate $\pi(m)\le m^2-1$ is an upper bound;
it does not assert that $\pi(m)$ divides $m^2-1$.
For example, direct iteration gives $\pi(5)=20$, and $20\ndivides24$.

A structural computation uses the roots of $X^2-X-1$.
Their $\sqrt5$ exists in $\F_p$ for $p=11$, but not for $p=7$.
Appendix~\ref{app:finite-fields} first builds arithmetic in an extension
field, then uses Frobenius and multiplicative orders to compute exact
periods. Those tools belong to a later unit; pair iteration and CRT
require only the arithmetic developed here.


\begin{exercisebox}
\begin{enumerate}
  \item Find the inverse of $37$ modulo $101$ using extended Euclid. \exlevel{Fluency}
  \item Solve $18x\congr30\pmod{42}$ and list all solutions modulo $42$. \exlevel{Fluency}
  \item Prove that congruence modulo $n$ is an equivalence relation. \exlevel{Reasoning}
  \item Show that an integer is divisible by $9$ exactly when its decimal digit sum is divisible by $9$. \exlevel{Reasoning}
  \item Solve $x\congr2\pmod7$, $x\congr3\pmod8$, $x\congr4\pmod9$ by the constructive CRT proof. \exlevel{Reasoning}
  \item Determine every possible last two digits of a square. Explain how CRT can split the work into moduli $4$ and $25$. \exlevel{Exploration}
  \item Prove that $x^2\congr1\pmod p$ has only the solutions $x\congr\pm1$ when $p$ is prime. Find a composite modulus where there are more than two solutions. \exlevel{Challenge}
  \item Run the forward extended-gcd table for $83$ and $29$, recording
  $r=83s+29t$ in every row. Find $29^{-1}\pmod{83}$ and check it by
  multiplication. \exlevel{Fluency}
  \item Solve $18x\congr24\pmod{42}$ using a B\'ezout certificate and
  list all residues modulo $42$. Explain why $18x\congr25\pmod{42}$
  has no solution. \exlevel{Fluency}
  \item For the CRT selectors, prove $e_i^2\congr e_i$,
  $e_ie_j\congr0$ when $i\ne j$, and $\sum_i e_i\congr1\pmod N$.
  Explain how these identities recover componentwise arithmetic.
  \exlevel{Reasoning}
  \item Solve $x\congr2\pmod6$, $x\congr5\pmod9$. Give the modulus
  of uniqueness. Why is $x\congr1\pmod6$, $x\congr2\pmod9$
  incompatible? \exlevel{Reasoning}
  \item Construct the polynomial of degree at most $2$ over $\Q$ with
  $f(-1)=2$, $f(0)=1$, $f(2)=3$ using Lagrange selectors.
  State all polynomials with these values. \exlevel{Fluency}
  \item Work in $\F_7$. Use Lagrange interpolation to find the
  polynomial of degree at most $2$ with $f(0)=1$, $f(1)=4$, $f(3)=2$.
  Certify each denominator inverse by an integer B\'ezout identity.
  \exlevel{Reasoning}
  \item Prove that no polynomial over $\Z/8\Z$ can satisfy
  $f(0)=0$, $f(2)=1$. Which hypothesis of the interpolation theorem
  fails? \exlevel{Reasoning}
  \item List the Fibonacci pair orbits modulo $2$ and $3$. Why is the
  first positive zero of $F_n\pmod3$ not the period? \exlevel{Fluency}
  \item Use CRT and $\pi(2)=3$, $\pi(3)=8$, $\pi(5)=20$ to find
  $\pi(6)$ and $\pi(15)$. Explain why the least common multiple,
  rather than the product, is the appropriate operation.
  \exlevel{Reasoning}
  \item Iterate $(a,b)\mapsto(b,2a+b)$ modulo $2$, starting at $(1,0)$.
  Show that the orbit has an initial segment before it repeats.
  Locate the failure of the Fibonacci reversibility proof.
  \exlevel{Reasoning}
  \item List the squares modulo $7$ and show that $X^2-X-1$ has no
  root in $\F_7$. Explain why the real Binet formula cannot simply
  be reduced modulo $7$ without first enlarging the arithmetic system.
  \exlevel{Bridge}
\end{enumerate}
\end{exercisebox}
